\documentclass[../../../../main.tex]{subfiles}
\begin{document}

\begin{flushright}
\footnotesize\textit{【自動文字起こし・要確認】}
\end{flushright}

{\bf [解]}
\begin{enumerate}
    \item[(1)] $y = e^{mx}$ の時, $y' = my$ だから,
    \begin{align*}
    3y' - 2y = 0 \iff (3m-2)y = 0 \iff m = \frac{2}{3} \quad (\because y \neq 0).
    \end{align*}

    \item[(2)] $y = e^{\frac{2}{3}x} \cdot u(x)$ と表す. $y' = e^{\frac{2}{3}x} \left(u'(x) + \frac{2}{3}u(x)\right)$ だから, ①に代入して
    \begin{align*}
    3 e^{\frac{2}{3}x} \left(u'(x) + \frac{2}{3}u(x)\right) - 2 e^{\frac{2}{3}x} u(x) = e^x
    \end{align*}
    \begin{align*}
    3 u'(x) = e^{\frac{1}{3}x}
    \end{align*}
    両辺積分して,
    \begin{align*}
    u(x) = e^{\frac{1}{3}x} + C.
    \end{align*}

    \item[(3)] (2)で $y = e^x + C \cdot e^{\frac{2}{3}x}$ で, $x=0$ で $y=10$ となるようにして,
    \begin{align*}
    10 = 1 + C \quad \therefore C = 9.
    \end{align*}
    だから,
    \begin{align*}
    y = e^x + 9 e^{\frac{2}{3}x}.
    \end{align*}
\end{enumerate}


\end{document}
